\subsection{Inventario exhaustivo de la superficie pública}
\label{sec:api-publica-exhaustiva}

La superficie canónica importable desde \codepath{hidden_attractors} contiene
99 símbolos. La tabla siguiente los enumera todos, sin convertir constantes,
tipos de datos o envolturas de flujo en procedimientos matemáticos que no son.
Las 33 funciones públicas de \codepath{hidden_attractors.plotting} se
documentan por separado, una por una, en la
Sección~\ref{sec:catalogo-graficas}. Las rutas especializadas de análisis e
integración que no se reexportan en la raíz se describen en sus secciones
temáticas del presente informe. El módulo experimental
\codepath{hidden_attractors.fractional} se documenta función por función en la
Sección~\ref{sec:nucleo-numerico-fraccionario}; entropía de cuencas, delay y
recurrencia avanzados se mantienen como imports module-qualified y se describen
en las Secciones~\ref{sec:basin-entropy-uncertainty},
\ref{sec:delay-embedding-advanced} y \ref{sec:recurrence-advanced}.

\begin{center}
\footnotesize
\begin{longtable}{@{}L{0.16\textwidth}L{0.50\textwidth}L{0.28\textwidth}@{}}
\toprule
Grupo & Símbolos públicos reexportados & Papel y fundamento \\
\midrule
\endfirsthead
\toprule
Grupo & Símbolos públicos reexportados & Papel y fundamento \\
\midrule
\endhead
Estabilidad de API &
\codepath{EXPERIMENTAL},
\codepath{INTERNAL},
\codepath{LEGACY},
\codepath{STABLE},
\codepath{api_tier},
\codepath{assert_tier},
\codepath{get_tier},
\codepath{PUBLIC_API_STABLE},
\codepath{PUBLIC_API_EXPERIMENTAL},
\codepath{PUBLIC_API_TIERS}
&
Metadatos y comprobaciones del contrato de estabilidad. No calculan una
propiedad dinámica. \\
\addlinespace

Modelo Chua &
\codepath{ChuaParameters},
\codepath{chua_parameters},
\codepath{chua_arctan_wu2023_parameters},
\codepath{chua_nonsmooth_parameters},
\codepath{equilibria_arctan},
\codepath{equilibria_nonsmooth},
\codepath{jacobian_arctan},
\codepath{jacobian_nonsmooth},
\codepath{rhs_arctan},
\codepath{rhs_nonsmooth}
&
Parámetros, campos \(F(X)\), equilibrios \(F(E)=0\) y Jacobianos
\(J=\partial F/\partial X\) de los modelos definidos en las Secciones 5 y 6. \\
\addlinespace

Registro de sistemas &
\codepath{ChaoticSystem},
\codepath{LureSystem},
\codepath{check_system_capability},
\codepath{get_system},
\codepath{known_workflows},
\codepath{list_systems},
\codepath{register_system},
\codepath{requirements_for}
&
Representación reusable del campo autónomo y, cuando aplica, del
desdoblamiento \(PX+b\psi(r^\top X)\); además valida capacidades exigidas por
cada flujo. \\
\addlinespace

Clases de destino e I/O &
\codepath{CLASS_LABELS},
\codepath{TARGET_CLASS_IDS},
\codepath{class_label},
\codepath{is_target_class},
\codepath{load_trajectory_csv}
&
Convenciones finitas de clasificación y carga portable de trayectorias. Una
clase de destino es una etiqueta bajo el contrato declarado, no una prueba
global de cuenca. \\
\addlinespace

Análisis y resultados &
\codepath{BifurcationPoint},
\codepath{LyapunovComputationRequest},
\codepath{LyapunovComputationSummary},
\codepath{LyapunovResult},
\codepath{PoincareCrossingResult},
\codepath{RobustnessCase},
\codepath{SpectrumResult},
\codepath{TimeSeriesLyapunovResult},
\codepath{bifurcation_points_from_trajectories},
\codepath{bifurcation_summary},
\codepath{compute_boundedness_metrics},
\codepath{compute_fft_psd},
\codepath{compute_lyapunov_spectrum},
\codepath{compute_trajectory_metrics},
\codepath{detect_poincare_crossings},
\codepath{estimate_time_series_lyapunov},
\codepath{integer_qr_benettin_lyapunov_exponents},
\codepath{integer_system_lyapunov_exponents},
\codepath{kaplan_yorke_dimension},
\codepath{trajectory_metrics},
\codepath{trajectory_metrics_for_system},
\codepath{validate_lyapunov_method_request},
\codepath{zero_one_test}
&
Tipos estructurados y funciones de caracterización independientes. Sus
ecuaciones, discretizaciones, unidades, llamadas y límites se presentan en la
Sección 12. \\
\addlinespace

Semillas armónicas &
\codepath{HarmonicSeed},
\codepath{find_harmonic_seed},
\codepath{find_lure_harmonic_seed},
\codepath{find_lure_omega_gain_candidates},
\codepath{find_omega_gain_candidates},
\codepath{validate_fractional_order}
&
Solución del cierre de balance armónico y validación del orden. Son
constructores de semillas; no certifican por sí mismos caos ni ocultedad. \\
\addlinespace

Contratos de flujo &
\codepath{BasinSliceSpec},
\codepath{DestinationClassifierSpec},
\codepath{FullWorkflowContract},
\codepath{ContinuationPlan},
\codepath{ContinuationTrace},
\codepath{DynamicReference},
\codepath{FINAL_LABELS},
\codepath{HiddennessTestResult},
\codepath{IntegratorSpec},
\codepath{NumericalContract},
\codepath{OFFICIAL_STAGE_ORDER},
\codepath{ParameterSweepSpec},
\codepath{RobustnessCaseSpec},
\codepath{RobustnessVerdict},
\codepath{PROTOCOL_VERSION},
\codepath{PostContinuationDecision},
\codepath{SEED_FAMILIES},
\codepath{SoftPrecheckResult},
\codepath{SphereControlSpec},
\codepath{StageEnvelope},
\codepath{StrictRefinementSpec},
\codepath{TargetReferenceSpec},
\codepath{TrajectoryDiagnosticsSpec},
\codepath{WorkflowInputSpec},
\codepath{UnifiedSeedRecord}
&
Tipos, enumeraciones y envolturas que fijan entradas, solver, memoria,
continuación, muestreo, evidencia y salidas. Formalizan el protocolo, pero no
añaden una ecuación dinámica distinta. \\
\addlinespace

Ejecución de flujos &
\codepath{continue_integer_lure_seed},
\codepath{final_integer_lure_attractor},
\codepath{integer_lure_seed},
\codepath{integrate_integer_lure},
\codepath{run_integer_lure_hiddenness_controls},
\codepath{validate_full_workflow_system},
\codepath{load_config},
\codepath{save_effective_config},
\codepath{run_attractor_only_workflow},
\codepath{run_bifurcation_workflow},
\codepath{run_basin_workflow},
\codepath{run_simple_workflow}
&
Orquestación reproducible de los métodos anteriores. Las cuatro últimas rutas
permiten simular, analizar o construir bifurcaciones y cuencas sin iniciar una
búsqueda de atractor oculto. \\
\bottomrule
\end{longtable}
\end{center}

El inventario anterior se limita de forma deliberada a la API soportada. Los
auxiliares privados, scripts de validación y backends internos permanecen
auditables en el código y en \codepath{docs/api_reference.md}, pero no se
promueven como llamadas de usuario. La documentación de un símbolo experimental
module-qualified tampoco implica promoción a la raíz ni a la categoría
\STABLE; hace visible su contrato para evitar llamadas ambiguas.
